% r03_r1_proof: the cumulant bound for the step Gaussian, part V: the thresholds, the small-volume case and the assembly over the four pavements. Input inside the proof of Theorem rot:thm:R1 in sec/r03_fill_K5a.tex.
\emph{The thresholds.} The threshold $b_*$ is the one fixed before \Cref{rot:lem:R1stg}; it depends only on $(n,\bar n,D,\gamma_0,c_P,\delta_P,c_K)$.

\emph{Small volume.} If some pavement consists of a single cell, then $|\Lambda|\le v_b$ by (G4) and no pavement is needed: with $\nu=\mu_P$ (centred), $\epsilon\defeq\mu_P(1-\chi_b)\le\epsilon_b$ by \Cref{rot:lem:R1tail} and $M\defeq\mathsf A|\Lambda|(c_Kb)^D\le\mathsf m_*$, \Cref{rot:lem:R1ext} moves the cutoff into the exponent at a cost of at most $2e^{2M}\epsilon$, \Cref{rot:lem:R1cum} expands to order $\bar n$ at a cost of $(2M)^{\bar n+1}$, and \Cref{rot:lem:R1chi} replaces $\kappa_{\mu_P}(V\chi_b;k)$ by $\kappa_{\mu_P}(V;k)$ at a cost of $k\,k!\,4^k(\mathsf A|\Lambda|)^k(2b)^{kD}\epsilon^{1/2}/k!$ for each $k$; the total is at most $E_*\le|J|E_*$ in both directions. (If $J=\emptyset$ then $V=0$ and the statement reduces to $0\ge\log\mu_P(\chi_b)\ge-4n|\Lambda|e^{-\gamma_0b^2/(2n)}$, which is \Cref{rot:lem:R1gci}.)

\emph{Assembly.} Otherwise iterate \Cref{rot:lem:R1stg} over $i=0,1,2,3$ with $J_0\defeq J$, $J_{i+1}\defeq J_i\cap\{\lambda_i<2.5w_b\}$ (the set $J''$ of stage $i$), and the thresholds $\beta_i\defeq\gamma_K^ib$ for the lower bound and $\beta_i\defeq\gamma_K^{-i}b$ for the upper bound, all of which lie in $[b_{\min},b_{\max}]$ together with $\beta_3/\gamma_K$. By (G3), $J_4\subset J\cap\bigcap_i\{\lambda_i<2.5w_b\}=\emptyset$, so $V_{J_4}=0$ and the cumulant increments of the four stages telescope to $\sum_{k\le\bar n}\kappa_{\mu_P}(V_J;k)/k!$; the numbers of $J$-cells add up to $\sum_iN_i\le4|J|$. The terminal factor is $\Phi(\gamma_K^4b;\emptyset)=\mu_P(\chi_{b_{\min}})\ge\exp\big(-4n|\Lambda|e^{-\gamma_0\gamma_K^8b^2/(2n)}\big)$ in the lower bound, by \Cref{rot:lem:R1gci} ($\mu_P$ is centred with component variances at most $\|C\|\le\gamma_0^{-1}$), and $\Phi(\gamma_K^{-4}b;\emptyset)=\mu_P(\chi_{b_{\max}})\le1$ in the upper bound. Hence
\[
\log\int\chi_be^Vd\mu_P-\sum_{k=1}^{\bar n}\frac{\kappa_{\mu_P}(V;k)}{k!}\ \le\ 4|J|E_*(b,\mathsf A),\qquad\log\int\chi_be^Vd\mu_P-\sum_{k=1}^{\bar n}\frac{\kappa_{\mu_P}(V;k)}{k!}\ \ge\ -4|J|E_*(b,\mathsf A)-4n|\Lambda|e^{-\gamma_0\gamma_K^8b^2/(2n)}.
\]

\emph{The constants.} $e_{\mathrm{cum}}=2(2M)^{\bar n+1}=2(2\mathsf m_*)^{\bar n+1}(\mathsf Ab^{D+6})^{\bar n+1}$ exactly. Every other term of $E_*$ is increasing in $\mathsf A$, so under $\mathsf Ab^{D+6}\le1$ it is at most its value at $\mathsf A=b^{-D-6}$; call the sum of these values $Q(b)$. It is a finite sum of polynomials in $b$ times one of $e^{-\varsigma w_b/2}$, $e^{-\delta_Pw_b}$, $e^{-\bar\mu w_b/4}$, $e^{-\mu_1w_b/8}$, $e^{-2\bar\mu w_b}$ and $\epsilon_{b_{\min}}^{1/2}$, each of which is at most a polynomial in $b$ times $e^{-\bar\mu b^{3/2}/8}$ (since $w_b\ge b^{3/2}$, $\bar\mu\le\mu_1\le\delta_P/2$ and $\bar\mu\le\varsigma$). With $\rho_3\defeq\bar\mu/16$ and
\[
S\defeq4\max\Big\{2(2\mathsf m_*)^{\bar n+1},\ \sup_{b\ge b_*}e^{\rho_3b^{3/2}}Q(b)\Big\}<\infty
\]
one has $4|J|E_*(b,\mathsf A)\le S|J|\big[(\mathsf Ab^{D+6})^{\bar n+1}+e^{-\rho_3b^{3/2}}\big]$, which is the statement. $S$, $b_*$, $\rho_3$ and $\gamma_K$ depend only on $(n,\bar n,D,\varsigma,\gamma_0,c_P,\delta_P,c_K)$; nothing depends on $\Lambda$, on $J$, on $\mathbb{L}$ or its period lattice, on $P$ and $V$ beyond the stated bounds, or on any translation invariance, of which none is used; no lower bound on $c_P$ and no upper bound on $P$ enters.
